\subsection{Algorithm Overview}

Adaptive recursive bisection replaces predetermined tree structures with data-driven selection at each split. The key principle: \textbf{at every binary split, evaluate both possible orderings and choose the one producing higher minority concentration}.

\textbf{Input}:
\begin{itemize}
\item Graph $G = (V, E)$ with vertex weights $w_{pop}(v)$ and edge weights $w_{ij}$
\item Target district count $k$
\item Minority percentage per vertex $m_v$ for $v \in V$
\item VRA threshold $\tau$ (e.g., 0.50 for majority-minority districts)
\end{itemize}

\textbf{Output}:
\begin{itemize}
\item Partition $\{P_1, \ldots, P_k\}$ with population balance, contiguity, and optimized minority concentration
\item Binary tree structure showing split sequence
\end{itemize}

\subsection{Pseudocode}

\begin{verbatim}
function AdaptiveBisection(G, k, minority_data):
    if k == 1:
        return {G}  # Base case: single district

    best_split = None
    best_score = -infinity
    best_ordering = None

    # Try all possible split sizes
    for k_L in range(1, k):
        k_R = k - k_L

        # Try both orderings: (k_L, k_R) and (k_R, k_L)
        for ordering in [(k_L, k_R), (k_R, k_L)]:
            # Perform binary split using METIS
            (G_L, G_R) = MetisBisect(G, ordering, minority_data)

            # Evaluate minority concentration
            score = EvaluateMinorityConcentration(G_L, G_R, minority_data)

            # Update best split if this ordering is better
            if score > best_score:
                best_score = score
                best_split = (G_L, G_R)
                best_ordering = ordering

    # Recurse on both subgraphs using locally optimal split
    (G_L, G_R) = best_split
    (k_L, k_R) = best_ordering

    districts_L = AdaptiveBisection(G_L, k_L, minority_data)
    districts_R = AdaptiveBisection(G_R, k_R, minority_data)

    return districts_L ∪ districts_R


function MetisBisect(G, (k_L, k_R), minority_data):
    # Compute target population for each subgraph
    W = TotalPopulation(G)
    target_L = W * (k_L / (k_L + k_R))
    target_R = W * (k_R / (k_L + k_R))

    # Call METIS with edge-weighted optimization
    # Edges between high-minority tracts have weight α > 1
    partition = METIS_PartGraphRecursive(
        nvtxs=|V|,
        xadj=adjacency_list,
        adjncy=neighbors,
        vwgt=vertex_weights,
        adjwgt=edge_weights,  # VRA-optimized weights
        nparts=2,
        tpwgts=[target_L/W, target_R/W],
        ubvec=[1.005]  # ±0.5% population tolerance
    )

    return (G[partition == 0], G[partition == 1])


function EvaluateMinorityConcentration(G_L, G_R, minority_data):
    # Compute maximum minority percentage achievable in subgraphs
    max_L = MaxMinorityInSubgraph(G_L, minority_data)
    max_R = MaxMinorityInSubgraph(G_R, minority_data)

    # Scoring function: sum of achievable MM districts
    score = 0
    if max_L >= 0.50:  # Left subgraph can create MM district
        score += 1
    if max_R >= 0.50:  # Right subgraph can create MM district
        score += 1

    # Tie-breaker: use actual maximum percentages
    score += 0.01 * (max_L + max_R)

    return score


function MaxMinorityInSubgraph(G, minority_data):
    # Upper bound: concentrate all minority population in one district
    total_minority = Σ_{v in G} m_v * w_{pop}(v)
    total_pop = Σ_{v in G} w_{pop}(v)
    target_district_pop = total_pop / k

    return min(total_minority / target_district_pop, 1.0)
\end{verbatim}

\subsection{Algorithm Details}

\subsubsection{Split Evaluation}

For each candidate split $(k_L, k_R)$, the algorithm performs METIS bisection with edge-weighted VRA optimization. METIS takes 2-4 seconds per split for typical state graphs (1000-5000 census tracts), making exhaustive evaluation feasible.

The evaluation function computes \textbf{maximum achievable minority percentage} in each subgraph:
\begin{equation}
m_{max}(G_i) = \min\left(\frac{\sum_{v \in V_i} m_v \cdot w_{pop}(v)}{|V_i| / k_i \cdot W}, 1.0\right)
\end{equation}

This upper bound assumes perfect concentration of all minority population in the subgraph into a single district. While actual partitioning may not achieve this theoretical maximum (due 	o geographic dispersion), it provides consistent ranking of candidate splits.

\subsubsection{Scoring Function}

The scoring function prioritizes splits that enable MM district creation:
\begin{equation}
\text{score}(G_L, G_R) = \mathbb{1}[m_{max}(G_L) \geq \tau] + \mathbb{1}[m_{max}(G_R) \geq \tau] + 0.01 \cdot (m_{max}(G_L) + m_{max}(G_R))
\end{equation}

The first two terms count potential MM districts (discrete). The third term breaks ties using actual minority percentages (continuous). This ensures splits enabling two MM districts dominate splits enabling only one, while splits with equal MM potential are ranked by concentration magnitude.

\subsubsection{Ordering Evaluation}

For each split size $k_L$, the algorithm evaluates both orderings $(k_L, k_R)$ and $(k_R, k_L)$ where $k_R = k - k_L$. This matters because METIS bisection is asymmetric: the algorithm may assign differently sized subgraphs 	o ``left'' vs ``right'' depending on initialization and refinement decisions.

\textbf{Example}: Splitting Alabama (k=7) into [2,5]:
\begin{itemize}
\item Ordering (2, 5): May assign northern counties 	o 2-district subgraph, southern Black belt 	o 5-district subgraph
\item Ordering (5, 2): May reverse this assignment
\end{itemize}

If Black belt concentration is higher, ordering (5, 2) likely produces better minority concentration in the smaller subgraph, affecting final MM district count.

\subsection{Complexity Analysis}

\textbf{Predetermined recursive bisection}: $(k-1)$ METIS bisections (one per split), each $O(|E|)$ for graph with $|E|$ edges. Total: $O(k \cdot |E|)$.

\textbf{Adaptive recursive bisection}: At each split level $\ell$, evaluates $O(k)$ candidate splits (split sizes 1 through $k-1$) with 2 orderings each. Number of split levels is $O(\log k)$ for balanced trees. Total: $O(k^2 \log k \cdot |E|)$ worst case.

\textbf{Practical performance}: For congressional redistricting with $k \leq 52$, adaptive bisection requires $\sim 200$ METIS calls vs $\sim 50$ for predetermined. With 2-second METIS runtime per call, adaptive takes 6-8 minutes vs 1-2 minutes. This overhead is negligible compared 	o data preparation (30+ minutes) and remains far faster than ensemble methods (hours 	o days).

\textbf{Optimization opportunity}: Prune candidate splits that cannot achieve VRA targets:
\begin{equation}
\text{if } m_{max}(G) < \tau \text{ for all subgraphs, skip split evaluation}
\end{equation}

This reduces adaptive bisection 	o $O(k \cdot |E|)$ in states where VRA compliance is infeasible (very low minority percentage or extreme geographic dispersion).

\subsection{Comparison 	o Related Algorithms}

\subsubsection{Exhaustive Tree Search}

The optimal tree structure could be found by enumerating all $C_{k-1}$ Catalan number distinct binary trees ($\approx 4^k / k^{3/2}$) and evaluating each. For $k=7$, this requires testing 42 trees; for $k=14$, 2,674,440 trees. Exhaustive search becomes infeasible for $k > 10$.

Adaptive bisection approximates exhaustive search through greedy local decisions, reducing complexity 	o polynomial while capturing most performance gain.

\subsubsection{Simulated Annealing on Trees}

An alternative approach: start with random tree, iteratively swap subtrees 	o improve objective. This requires defining neighborhood (tree edit distance) and running multiple iterations (typically 1000+). Computational cost similar 	o adaptive bisection but may find better structures.

We do not pursue this approach because: (1) loses determinism (random initialization), (2) requires tuning annealing schedule, (3) adaptive bisection already closes most gap between predetermined and n-way (see Results section).

\subsubsection{Look-Ahead Search}

Instead of evaluating only immediate split, look ahead multiple levels (e.g., evaluate complete 2-level tree structure). This increases search space exponentially: 2-level look-ahead requires $O(k^4)$ evaluations vs $O(k^2)$ for greedy.

We implement 1-level (greedy) adaptive bisection as baseline. Future work may investigate k-level look-ahead for small states ($k < 10$).

\subsection{Implementation}

Open-source implementation available at \texttt{https://github.com/[anonymous]}.

\textbf{Key components}:
\begin{itemize}
\item \texttt{adaptive\_bisection.py}: Main algorithm with split evaluation
\item \texttt{metis\_wrapper.py}: Python interface 	o METIS library
\item \texttt{scoring.py}: Minority concentration evaluation functions
\item \texttt{visualization.py}: Tree structure and district map generation
\end{itemize}

\textbf{Dependencies}:
\begin{itemize}
\item METIS 5.1.0+ (graph partitioning)
\item Python 3.8+ with NumPy, pandas, GeoPandas
\item Census tract data (from Census Bureau)
\end{itemize}

\textbf{Runtime}: Alabama (k=7, 1,388 census tracts): 8 minutes for adaptive bisection, 2 minutes for predetermined tree. Georgia (k=14, 3,247 tracts): 22 minutes adaptive, 5 minutes predetermined.

\subsection{Determinism and Reproducibility}

METIS uses randomized algorithms for initial partitioning and refinement. To ensure reproducibility, fix random seed via \texttt{METIS\_option[METIS\_OPTION\_SEED] = 42}.

With fixed seed, adaptive bisection produces identical results across runs. This enables:
\begin{itemize}
\item Verification by independent parties
\item Legal defensibility (courts can reproduce exact districts)
\item Comparison across multiple states and census years
\end{itemize}

Alternative: run multiple seeds (e.g., 10 runs with seeds 1-10) and select median result. This provides robustness 	o METIS randomness while maintaining computational feasibility.
